\subsubsection{Reelle Zahlen}
\label{subsubsec:fundamentals_real_numbers}

Zahlen, die sich nicht als Quotient ganzer Zahlen darstellen lassen, heißen
\emph{irrationale Zahlen}. Beispiele sind
\[
  \sqrt2=1{,}4142\ldots,
  \qquad
  \pi=3{,}14159\ldots,
  \qquad
  \mathrm e=2{,}71828\ldots.
\]
Rationale und irrationale Zahlen bilden zusammen die \emph{reellen Zahlen}.

\begin{center}
  \begingroup
    \def\NumberSetStage{4}
    \CachedTikzFigure{content/resources/chapters/fundamentals/tikz/number_set_inclusions.tex}
  \endgroup
\end{center}

Die reellen Zahlen $\R$ sind bezüglich Addition, Subtraktion und Multiplikation
sowie bezüglich der Division durch eine von null verschiedene Zahl
abgeschlossen.

\begin{KeyPointBox}{Vollständigkeit der reellen Zahlen}{fundamentals_real_number_completeness}
  Jeder Punkt auf der Zahlengeraden entspricht genau einer reellen Zahl und
  umgekehrt. Anschaulich besitzt die Zahlengerade daher keine Lücken: Die
  reellen Zahlen $\R$ sind \emph{vollständig}.
  \begin{center}
    \CachedTikzFigure{content/resources/chapters/fundamentals/tikz/real_number_line.tex}
  \end{center}
   Damit bilden sie eine zentrale Zahlenmenge in den Naturwissenschaften und der Technik.
\end{KeyPointBox}

Bezüglich des Wurzelziehens sind allerdings auch die reellen Zahlen nicht abgeschlossen:
Die reelle Gleichung
\[
  x^2=-1
\]
besitzt in $\R$ keine Lösung, denn das Quadrat jeder reellen Zahl ist
nichtnegativ.


